
Our experiments test the hypothesis: Do data quality components Q(D) correlate with information density?

\subsubsection{Datasets}

We created 12 controlled C4 subsets (10GB each, 120GB total) using fractional factorial sampling. Each subset isolates specific quality variations while holding other factors approximately constant via stratified sampling. Document selection uses MinHash LSH for deduplication detection (128 hash functions, Jaccard threshold 0.9), GPT-2 small perplexity scoring, automated domain classification, and efficiency filtering.

\subsubsection{Metrics}

For each subset, we compute:

\textbf{Q(D) components:}
\begin{itemize}
\item Deduplication ratio: 1 - (duplicate\_tokens / total\_tokens)
\item Domain diversity: Shannon entropy over 8 C4 source categories
\item Perplexity score: 1 - normalized log(GPT-2\_perplexity)
\item Token efficiency: content\_tokens / total\_tokens
\end{itemize}

\textbf{Information density:}
\begin{itemize}
\item Token entropy: H(unigram distribution) normalized by $log_{2}|V|$
\item N-gram redundancy: 1 - (repeated 3-grams / total 3-grams)
\item Semantic diversity: 1 - average cosine similarity of SentenceTransformer embeddings
\end{itemize}

Combined info\_density = average of three normalized components $\in$ [0,1].

\subsubsection{Evaluation Protocol}

\begin{enumerate}
\item Correlation analysis: Compute Pearson r and Spearman $\rho$ between each Q(D) component and info\_density across 12 subsets
\item Statistical significance: Two-tailed t-test for $H_0$: r = 0 at $\alpha$ = 0.01 with Bonferroni correction
\item Effect size: Report $R^2$ (variance explained) for composite Q(D) fitted via linear regression
\item Success criterion: ALL four components must achieve r > 0.5, p < 0.01
\end{enumerate}

\subsubsection{Baseline Comparisons}

\begin{itemize}
\item Random baseline: Correlate Q(D) with shuffled info\_density
\item Single-metric baseline: Best individual component
\item Composite advantage: Does weighted Q(D) outperform best single metric?
\end{itemize}

\subsubsection{Compute Resources}

\begin{itemize}
\item V100 GPUs (32GB), 9.2 GPU-hours total for all subset creation and analysis
\item Environment: PyTorch 2.6.0, Transformers 4.46.0, datasketch 1.6.0
\end{itemize}
